% ============================================================
\section{Funtor \texorpdfstring{$H: End \to Graph$}{H: End -> Graph}}
\label{sec:functor-end-graph}

\subsection{Descrição das Categorias}

\begin{definition}[Categoria $\mathbf{End}$ (modelo MATLAB/Octave)]
Os objetos são pares $(X, \varphi)$ onde $X$ é um conjunto finito e $\varphi: X \to X$ é uma endoaplicação. Em MATLAB, representa-se $X = \{1, \ldots, n\}$ implicitamente pelo seu cardinal $n$, e $\varphi$ como um vetor de comprimento $n$ com entradas em $\{1, \ldots, n\}$. Um morfismo $f: (X_1, \varphi_1) \to (X_2, \varphi_2)$ é uma aplicação que comuta com as endoaplicações: $f \circ \varphi_1 = \varphi_2 \circ f$.
\end{definition}

\begin{lstlisting}[style=matlab, caption={Objeto e morfismo em $\mathbf{End}$}]
% Objeto: conjunto {1,2,3} com endoaplicacao phi(1)=2, phi(2)=3, phi(3)=1
phi = [2, 3, 1];

% Verificacao de morfismo f: (X1,phi1) -> (X2,phi2)
% f comuta com as endoaplicacoes se f(phi1(x)) == phi2(f(x)) para todo x
f    = [1, 2, 3];   % exemplo de morfismo candidato
ok   = all(f(phi1) == phi2(f));  % deve ser true
\end{lstlisting}

\begin{definition}[Categoria $\mathbf{Graph}$ (modelo MATLAB/Octave)]
Os objetos são grafos dirigidos descritos por um conjunto de vértices $V$, um conjunto de arestas $A$, e dois mapas $s, t: A \to V$ que indicam, respetivamente, a origem e o destino de cada aresta. Um morfismo de grafos é um par $(f_V, f_A)$ que preserva a estrutura de incidência: $f_V \circ s = s' \circ f_A$ e $f_V \circ t = t' \circ f_A$.
\end{definition}

\begin{lstlisting}[style=matlab, caption={Objeto e morfismo em $\mathbf{Graph}$}]
% Grafo dirigido com 3 vertices e 3 arestas
Vertices = [1, 2, 3];
Arestas  = [1, 2, 3];
Origem   = [1, 2, 3];   % s(a_i) = i
Destino  = [2, 3, 1];   % t(a_i) = phi(i)

% Verificacao de morfismo de grafos (fV, fA)
comuta_s = isequal(fV(Origem1(fA)), Origem2);   % fV o s1 = s2 o fA
comuta_t = isequal(fV(Destino1(fA)), Destino2); % fV o t1 = t2 o fA
\end{lstlisting}

\subsection{Funtor Covariante \texorpdfstring{$H: End \to Graph$}{H: End -> Graph}}

O funtor $H$ constrói canonicamente, a partir de qualquer endoaplicação $\varphi: X \to X$, um grafo dirigido em que cada elemento de $X$ é simultaneamente vértice e aresta, com a origem dada pela identidade e o destino dado por $\varphi$.

\begin{itemize}
    \item \textbf{Nos objetos:} Uma endoaplicação $\varphi$ de tamanho $n$ transforma-se num grafo onde:
    \begin{itemize}
        \item Os \textbf{vértices} ($V$) são os elementos $\{1, \ldots, n\}$;
        \item As \textbf{arestas} ($A$) são indexadas também por $\{1, \ldots, n\}$;
        \item A função \textbf{origem} ($s$) é a identidade: a aresta $i$ sai do vértice $i$;
        \item A função \textbf{destino} ($t$) é a aplicação $\varphi$: a aresta $i$ entra no vértice $\varphi(i)$.
    \end{itemize}
    \item \textbf{Nos morfismos:} Um morfismo $f: (X_1, \varphi_1) \to (X_2, \varphi_2)$ em $\mathbf{End}$ (que comuta com as endoaplicações) traduz-se num homomorfismo de grafos $(f_V, f_A) = (f, f)$ que mapeia vértices e arestas pelo mesmo mapa $f$, preservando a estrutura de incidência.
\end{itemize}

\begin{lstlisting}[style=matlab, caption={Funtor $H$ aplicado a uma endoaplicação $\varphi: \{1,2,3\} \to \{1,2,3\}$}]
% Endoaplicacao phi: ciclo fechado
phi = [2, 3, 1];

% Traducao para grafo dirigido
Vertices = 1:length(phi);
Arestas  = 1:length(phi);
Origem   = Vertices;   % s(i) = i  (identidade)
Destino  = phi;        % t(i) = phi(i)
% Resultado: grafo com ciclo dirigido 1->2->3->1
\end{lstlisting}

\begin{proposition}
O funtor $H: \mathbf{End} \to \mathbf{Graph}$ satisfaz as seguintes propriedades:
\begin{itemize}
    \item $H(\mathrm{id}_X)$ é o homomorfismo de grafos identidade;
    \item $H(g \circ f) = H(g) \circ H(f)$ (preservação da composição);
    \item Pontos fixos de $\varphi$ correspondem a laços no grafo resultante ($s(a) = t(a)$).
\end{itemize}
\end{proposition}

\subsection{Transformações Naturais \texorpdfstring{$\gamma: H_1 \Rightarrow H_2$}{gamma: H1 => H2}}

Dados dois funtores paralelos $H_1, H_2: \mathbf{End} \to \mathbf{Graph}$, uma transformação natural $\gamma: H_1 \Rightarrow H_2$ é uma família de homomorfismos de grafos,
\[
    \gamma_{(X,\varphi)} : H_1(X, \varphi) \to H_2(X, \varphi), \quad \text{para cada objeto } (X,\varphi) \in \mathbf{End},
\]
tal que para cada morfismo $f: (X_1, \varphi_1) \to (X_2, \varphi_2)$ o seguinte quadrado comuta:
\[
    H_2(f) \circ \gamma_{(X_1,\varphi_1)} \;=\; \gamma_{(X_2,\varphi_2)} \circ H_1(f).
\]
Considerando $H_1$ como o funtor de passo único (arestas $x \to \varphi(x)$) e $H_2$ como o funtor de passo duplo (arestas $\varphi(x) \to \varphi(\varphi(x))$), a transformação natural $\gamma: H_1 \Rightarrow H_2$ tem componentes dadas pelo homomorfismo de grafos onde os vértices avançam de acordo com $\varphi$ e as arestas se mantêm, comutando o diagrama de incidência.

\begin{lstlisting}[style=matlab, caption={Verificação da transformação natural $\gamma: H_1 \Rightarrow H_2$ em $\mathbf{End} \to \mathbf{Graph}$}]
X   = 1:4;
phi = [2, 3, 3, 1];

% Funtor H1: grafo de passo unico
Origem1  = X;       Destino1 = phi;

% Funtor H2: grafo de passo duplo
Origem2  = phi;     Destino2 = phi(phi);

% Componentes da transformacao natural gamma
gamma_V = phi;   % transformacao dos vertices: v -> phi(v)
gamma_E = X;     % transformacao das arestas: identidade nos indices

% Validacao da comutatividade dos diagramas de incidencia
comuta_origem  = isequal(gamma_V(Origem1),  Origem2(gamma_E));  % true
comuta_destino = isequal(gamma_V(Destino1), Destino2(gamma_E)); % true
\end{lstlisting}